% Generated draft: requires mathematical and source review.
\documentclass[11pt]{article}
\usepackage[margin=25mm]{geometry}
\usepackage{fontspec}
\setmainfont{Latin Modern Roman}
\usepackage{amsmath,amssymb,mathtools,unicode-math}
\setmathfont{Latin Modern Math}
\usepackage{enumitem,needspace,xurl,hyperref}
\setlength{\parindent}{0pt}
\setlength{\parskip}{5pt}
\setlength{\emergencystretch}{3em}
\newcommand{\N}{\mathbb{N}}
\newcommand{\Z}{\mathbb{Z}}
\newcommand{\Q}{\mathbb{Q}}
\newcommand{\R}{\mathbb{R}}
\newcommand{\C}{\mathbb{C}}
\newcommand{\F}{\mathbb{F}}
\newcommand{\A}{\mathbb{A}}
\newcommand{\PP}{\mathbb{P}}
\newcommand{\E}{\mathbb{E}}
\newcommand{\eps}{\varepsilon}
\newcommand{\dd}{\,\mathrm{d}}
\begin{document}
\section*{UnsolvedMath ID: 30001120}
\textbf{Source problem code:} catalogue entry
\subsection{Definitions and mathematical statement}
Let \(d\) denote a positive integer and let \(\Z^d\) be the integer lattice. Consider the parabolic Anderson model on this lattice with random spatial potential, denoted here by \(\xi\). The supplied record denotes its remaining-lifetime process by \(R(t)\), but gives no mathematical definition of that process.

\textbf{Source question.} At what tail threshold does the ageing behaviour described by \(R(t)\) change qualitatively when the distribution of the random potential has a tail lighter than the polynomial tails treated in the cited Oberwolfach report?

\textbf{Unresolved formulation.} The record does not specify the potential's distribution, dependence assumptions or admissible tail classes; the evolution equation, its normalization or initial condition; or the definition of \(R(t)\). It also does not define the comparison meant by ``lighter than polynomial'', the qualitative change to be detected, or a limiting distribution or time normalization that would characterize ageing. Thus the supplied question identifies a research target, but does not determine a unique quantified mathematical conjecture or establish that a single tail threshold exists. These missing specifications must be recovered from the cited source before a more precise statement can be given.
\subsection{Short English statement}
Determine where ageing in the parabolic Anderson model on \(\Z^d\), as measured by the remaining-lifetime process \(R(t)\), changes as potential tails become lighter than the polynomial tails considered in the source report. The supplied record leaves the process, tail comparison and criterion for change undefined, so this is an unresolved research-question formulation requiring human mathematical review.
\subsection{Sources}
\begin{itemize}
\item Supplied local UnsolvedMath registry record: permanent ID \(30001120\), title \emph{Tail-Dependent Ageing in the Parabolic Anderson Model}, exact source problem code \texttt{catalogue entry}, in the Oberwolfach Reports Open Problems collection. No individual proposer is supplied; the proposed year is recorded as 2008.
\item Supplied source citation: \emph{Interplay of Analysis and Probability in Physics} (2009), Oberwolfach Reports, pp.~3065--3138, DOI 10.4171/owr/2008/55. Original OWR report: \url{https://doi.org/10.4171/owr/2008/55}. The supplied background describes these reports as documenting mathematical workshops and their open-problem discussions.
\item Supplied literature reference labelled ``Weibull localization and ageing'': \url{https://doi.org/10.1214/13-AOP882}.
\item Supplied literature reference labelled ``polynomial-tail ageing'': \url{https://doi.org/10.1214/10-AIHP394}.
\item Literature assessment checked on \texttt{2026-08-21}: this is supplied dated curation, not independently verified here. It labels the problem partially solved and reports that ageing and localization are proved for polynomial, Weibull and doubly-exponential tails, including explicit residual-lifetime asymptotics in several classes, while a complete theory for all lighter tails and the qualitative transition remains open. The supplied background describes this assessment as dated literature triage derived from OWR curation and calls for independent verification of time-sensitive status and cited claims.
\item The registry labels the statement \texttt{corrected\_verified} and says its mathematical content was checked against the cited official Oberwolfach report, with notation, context and question form normalized without changing the source-level claim. This is supplied verification metadata, not an independent source check performed for this draft.
\end{itemize}
\end{document}
